\section{Design and Methodology}

\subsection{Proposed System}
The proposed system has three major modules (Figure~\ref{fig:flowextract} and Figure~\ref{fig:flowquery}):
\begin{enumerate}
\item \textbf{Extraction pipeline.} Converts any supported file into a common grid model, recognises
      sections, time rows, lab batches, legends and breaks, validates every entry and gives it a
      status: \textsc{Verified}, \textsc{Unverified}, \textsc{Incomplete} or \textsc{Rejected}.
\item \textbf{Back end.} Stores each timetable as an archive, enforces role-based access, detects
      clashes and answers questions through a deterministic \emph{compositional query planner}.
\item \textbf{Cross-platform app.} Role dashboards, weekly view of any batch, teacher or room, the
      Ask screen, archives, uploads, review screens and administrator tools.
\end{enumerate}
\begin{figure}[H]
\centering
\begin{tikzpicture}[node distance=4.5mm]
\node[term] (s) {Start: job claimed};
\node[blk,below=of s] (a) {Detect format (magic bytes)};
\node[dec,below=of a] (d1) {PDF?};
\node[blk,right=12mm of d1,text width=40mm] (o) {Office/CSV/image adapter\\(openpyxl, python-docx, OCR)};
\node[blk,below=of d1] (b) {Per page: text + words with coordinates};
\node[dec,below=of b] (d2) {enough\\text?};
\node[blk,right=12mm of d2] (ocr) {Tesseract OCR};
\node[blk,below=of d2] (c) {Ruled tables (keep merged spans)\\or geometric gutter fallback};
\node[blk,below=of c] (g) {Common Grid model};
\node[blk,below=of g] (h) {Sections, W.E.F., tentative; time rows (Alg.~\ref{alg:time})};
\node[blk,below=of h] (i) {Cells $\rightarrow$ entries (Sec.~4.4); legend expansion};
\node[blk,below=of i] (j) {Map labels to master data (unique exact match)};
\node[blk,below=of j] (k) {Validate, status per entry, timetable status};
\node[term,below=of k] (e) {Persist + notify};
\draw[arr] (s)--(a); \draw[arr] (a)--(d1); \draw[arr] (d1)--node[lbl,right]{yes}(b); \draw[arr] (d1)--node[lbl,above]{no}(o);
\draw[arr] (o.south) |- (g.east);
\draw[arr] (b)--(d2); \draw[arr] (d2)--node[lbl,right]{yes}(c); \draw[arr] (d2)--node[lbl,above]{no}(ocr);
\draw[arr] (ocr.south) |- (c.east);
\draw[arr] (c)--(g); \draw[arr] (g)--(h); \draw[arr] (h)--(i); \draw[arr] (i)--(j); \draw[arr] (j)--(k); \draw[arr] (k)--(e);
\end{tikzpicture}
\caption{Flowchart of the extraction pipeline}
\label{fig:flowextract}
\end{figure}


\subsection{Dataset}
The system was developed and evaluated on the Computer Engineering department's real timetable for
the current semester: a 10-page PDF covering S.E., T.E., B.E.\ and M.Tech.\ divisions with 453 classes
(570 rows including breaks and activities). Problems found in this document shaped the design: lab
cells with four batches, codes expanded only in each page's legend, the noon row printed
``12.15 a.m.\ to 01.15 p.m.'', a combined B.E.\ elective page, LLC slots with no teacher or room, and
spelling variants (``Sushama''/``Sushma'' Pande, ``Discreate'' Structures). A smaller reference PDF
and synthetic Excel, Word and scanned files were used for the other input formats.

\subsection{Time-Row Resolution by Constrained Backtracking}
Each time label is turned into a small set of candidate intervals. The algorithm searches for
assignments in which every row is plausible (between 7:00 and 21:00, at most 6 hours) and no row
starts before the previous row ends. A reading is accepted as certain only if it is the \emph{same in
every valid assignment}; otherwise the row goes to review. The college's noon row is therefore read as
12:15--13:15 only because the row before ends at 12:15 and the row after starts at 13:15. The search is
bounded at 64 assignments, so it costs $O(64r)$ for $r$ rows.
\begin{algorithm}[H]
\caption{Resolve a column of time labels}\label{alg:time}
\begin{algorithmic}[1]
\Require labels $L_1..L_r$ (top to bottom); each $L_i$ gives start/end hour, minute and optional marker
\Ensure for each row: (start, end, certain?) or a review flag
\For{$i \gets 1$ \textbf{to} $r$}
  \State $O_i \gets \{(s,e) : s \in \textsc{Cand}(L_i.\mathit{start}), e \in \textsc{Cand}(L_i.\mathit{end}), \textsc{Plausible}(s,e)\}$ \Comment{literal markers honoured}
  \If{$O_i = \emptyset$} $O_i \gets$ readings ignoring markers, flagged \textit{non-literal} \EndIf
\EndFor
\State $S \gets \emptyset$
\Procedure{Walk}{$i, \mathit{prevEnd}, A$}
  \If{$|S| > 64$} \Return \EndIf \Comment{bound the search}
  \If{$i > r$} $S \gets S \cup \{A\}$; \Return \EndIf
  \ForAll{$(s,e) \in O_i$ \textbf{with} $s \geq \mathit{prevEnd}$}
     \State \Call{Walk}{$i+1, e, A \cdot (s,e)$}
  \EndFor
\EndProcedure
\State \Call{Walk}{$1, 0, \langle\rangle$}
\For{$i \gets 1$ \textbf{to} $r$}
  \State $R_i \gets \{A_i : A \in S\}$
  \If{$|R_i| = 1$ \textbf{and} reading is literal} mark row $i$ \textsc{certain}
  \ElsIf{$|R_i| = 1$} keep reading, set \texttt{time\_uncertain}, add \texttt{TIME\_LABEL\_INCONSISTENT}
  \Else{} leave time empty, add \texttt{TIME\_AMBIGUOUS\_AMPM} (entries become \textsc{Incomplete})
  \EndIf
\EndFor
\end{algorithmic}
\end{algorithm}


\subsection{Cell Parsing}
A cell is split into lines, lines belonging to one lab batch (e.g.\ \texttt{A1 COA Lab JS 606-4}) are
grouped, and the batch (\texttt{[A-D][1-4]}), room, teacher codes and subject are identified. Codes are
expanded only if they appear in the legend of the \emph{same page}; otherwise the code is kept and an
informational flag is raised.

\subsection{Conflict Detection (Group, Sort and Sweep)}
Entries are grouped by resource (teacher, room, batch) and day, sorted by start time and swept forward;
the inner loop stops as soon as a later entry starts after the current one ends, giving
$O(n\log n + m)$ instead of $O(n^2)$~\cite{cormen2009}. Intended overlaps (the same source cell,
parallel lab batches, elective options, a combined session printed on several pages) are not
reported. The overlap test follows Allen's interval relations~\cite{allen1983}.
\begin{algorithm}[H]
\caption{Detect professor, room and batch clashes}\label{alg:conflict}
\begin{algorithmic}[1]
\Require entries $E$ of one timetable (classes with day and times, not rejected)
\Ensure list of conflicts with confidence \textsc{Confirmed}/\textsc{Possible}
\State $G \gets$ empty hash map
\ForAll{$x \in E$}
  \ForAll{resource $(k, \mathit{key}) \in \{$professor, room, batch$\}$ of $x$}
     \State append $x$ to $G[(k, \mathit{key}, x.\mathit{day})]$ \Comment{id if mapped, else normalised label}
  \EndFor
\EndFor
\ForAll{group $g \in G$}
   \State sort $g$ by start time
   \For{$i \gets 1$ \textbf{to} $|g|$}
     \For{$j \gets i+1$ \textbf{to} $|g|$}
        \If{$g_j.\mathit{start} \geq g_i.\mathit{end}$} \textbf{break} \Comment{sorted: no later entry can overlap}\EndIf
        \If{$g_i.\mathit{start} < g_j.\mathit{end} \wedge g_j.\mathit{start} < g_i.\mathit{end}$ \textbf{and not} \textsc{Intended}$(g_i, g_j)$}
            \State report $(k, g_i, g_j)$; \textsc{Confirmed} if both verified, mapped and certain
        \EndIf
     \EndFor
   \EndFor
\EndFor
\Function{Intended}{$a, b$}: same source cell, same combined session on several pages, parallel
elective options, or a sub-group's own written period \EndFunction
\end{algorithmic}
\end{algorithm}


\subsection{Compositional Query Planner}
The engine does not match a whole question to one of a fixed list. It recognises independent pieces
and combines them into a typed \textbf{query plan}:
\begin{itemize}
\item \textbf{Operation} (17): \textsc{List}, \textsc{Count}, \textsc{Exists}, \textsc{Total\_Hours},
      \textsc{First}, \textsc{Last}, \textsc{Next}, \textsc{Per\_Day}, \textsc{Span},
      \textsc{Free\_Slots}, \textsc{Breaks}, \textsc{After/Before\_Course}, \textsc{Teachers},
      \textsc{Per\_Room}, \textsc{Room\_Info}, \textsc{Conflicts}.
\item \textbf{Filters}: me, batch, professor, room, course; lab or lecture; room type.
\item \textbf{When}: a date, a weekday, this or next week, or the whole week.
\item \textbf{Time window}: at, after, before, between, after lunch, now.
\item \textbf{Metric and sort}: most/least classes or hours, earliest/latest start or finish.
\end{itemize}
For example, ``Which day do I finish earliest?'' becomes
\texttt{\{operation: PER\_DAY, metric: LAST\_END, sort: EARLIEST\}}. Names are matched against a
vocabulary built from the timetable itself: a teacher can be named by code (``AVN''), first name,
surname or ``Nimkar sir''; entries sharing a faculty code are merged; subjects match partial words and
one-letter typos. The plan is validated against a closed registry and compiled into pre-written
parameterised SQL, so user text never becomes SQL (safe against the attacks
in~\cite{greshake2023,pedro2023}). Ambiguous questions (``at 2'', a surname shared by two teachers)
get a clarification instead of a guess, and every answer carries its plan.
\begin{algorithm}[H]
\caption{Build a query plan from the parsed question}\label{alg:plan}
\begin{algorithmic}[1]
\Require parsed question $p$, specialised operation $S$ recognised by the rule cascade (or $\bot$)
\Ensure plan $P$
\If{$S$ is pinned (archives, primary, cross-domain, time-until)} \Return $P(S)$ \EndIf
\State $k \gets$ \textsc{Lab} / \textsc{Lecture} from words (``lab'', ``practical'' / ``lecture'', ``theory''); $\bot$ otherwise
\State $w \gets$ window from times and words (between, before, after, after lunch, at)
\State $o \gets$ \textsc{Total\_Hours} (``total hours'') $\vee$ \textsc{Count} (``how many'') $\vee$ \textsc{Per\_Day} (``which day'', ``busiest'') $\vee$ \textsc{Next} / \textsc{First} / \textsc{Last} $\vee$ \textsc{Exists} (``any lab'')
\If{$o = $ \textsc{Per\_Day}} choose metric (finish $\rightarrow$ \textsc{Last\_End}, start $\rightarrow$ \textsc{First\_Start}, else classes/hours) and sort (most/least, earliest/latest) \EndIf
\If{$o = $ \textsc{Count} and $S$ returns rooms/professors/courses} \Return $P(\textsc{Count}$ over $S)$ \EndIf
\If{$o = \bot$ and $k = \bot$ and $w \neq$ \textsc{Before}} \Return $P(S)$ \Comment{specialised operation unchanged} \EndIf
\State \Return composed $P(o \text{ or } \textsc{List}, k, w, \text{filters}, \text{when}, \text{sort})$
\end{algorithmic}
\end{algorithm}

\begin{figure}[H]
\centering
\begin{tikzpicture}[node distance=4.5mm]
\node[term] (s) {Question text};
\node[blk,below=of s] (p) {Parse: dates, times, entities};
\node[dec,below=of p] (x) {cross-\\domain?};
\node[blk,right=12mm of x,fill=red!10] (x1) {UNSUPPORTED (Q30)};
\node[blk,below=of x] (c) {Resolve context (server)};
\node[blk,below=of c] (v) {Scan vocabulary of this timetable};
\node[dec,below=of v] (k) {plan\\built?};
\node[blk,right=12mm of k,fill=red!10] (k1) {UNSUPPORTED +\\examples};
\node[dec,below=of k] (r) {allowed for\\role/domain?};
\node[blk,right=12mm of r,fill=red!10] (r1) {Refuse (reason)};
\node[dec,below=of r] (m) {params\\complete?};
\node[blk,right=12mm of m,fill=yellow!20] (m1) {CLARIFICATION\\(choices)};
\node[blk,below=of m] (t) {Run template QTnn (bound params, role scope)};
\node[term,below=of t] (e) {OK / NO\_MATCH / DATA\_UNVERIFIED};
\draw[arr] (s)--(p); \draw[arr] (p)--(x); \draw[arr] (x)--node[lbl,above]{yes}(x1); \draw[arr] (x)--node[lbl,right]{no}(c);
\draw[arr] (c)--(v); \draw[arr] (v)--(k); \draw[arr] (k)--node[lbl,above]{no}(k1); \draw[arr] (k)--node[lbl,right]{yes}(r);
\draw[arr] (r)--node[lbl,above]{no}(r1); \draw[arr] (r)--node[lbl,right]{yes}(m);
\draw[arr] (m)--node[lbl,above]{no}(m1); \draw[arr] (m)--node[lbl,right]{yes}(t); \draw[arr] (t)--(e);
\end{tikzpicture}
\caption{Flowchart of the query engine}
\label{fig:flowquery}
\end{figure}


\subsection{Security Design}
Passwords are hashed with Argon2id~\cite{biryukov2016}. Access tokens are JWTs~\cite{rfc7519} valid for
15 minutes; refresh tokens rotate, are stored only as hashes, and reuse of a revoked token revokes the
whole family. Roles are read from the database on every request, never from the client. Requests for
another user's personal timetable return 404 and are logged. Table~\ref{tab:security} lists the main
threats and controls.
\begin{table}[H]
\centering
\caption{Threats and controls}
\label{tab:security}
\small
\begin{tabularx}{\textwidth}{@{}p{3.6cm}X@{}}
\toprule
\textbf{Threat} & \textbf{Control in the design}\\
\midrule
Password theft & Argon2id hashing~\cite{biryukov2016}; generic login error; per-IP and per-e-mail rate limit\\
Stolen refresh token & Rotation on every refresh; reuse of an old token revokes the whole token family\\
Viewing policy & The official timetable is public, so every signed-in user may view and ask about all of it; roles only decide what ``my'' means and who may change data (Admin)\\
Role escalation & Roles read from \texttt{USER\_ROLE} on every request; role change bumps \texttt{token\_version}\\
IDOR on personal data & Owner check in the query itself; non-owners get 404 (no existence leak) and an \texttt{ACCESS\_DENIED} audit row\\
SQL / prompt injection & No SQL built from text; query plans use a closed vocabulary of operations and filters; typed, bound parameters; unknown parameters rejected\\
Malicious upload & Extension allow-list \emph{and} magic-byte signature check; size/page limits; private storage; parsing isolated in the worker; LibreOffice with isolated profile and timeout\\
Account takeover via reset & Reset codes never returned in API responses or logs in production (bug P0-1 found and fixed in our audit)\\
Information leakage & Uniform error envelope with request id; no stack traces or SQL in responses; no tokens or free text in logs\\
Secrets in the app & Only the public API URL is in the bundle; tokens in iOS Keychain / Android Keystore (SecureStore)\\
\bottomrule
\end{tabularx}
\end{table}

\newpage
